\documentclass[11pt]{article}
% Standalone preamble for web-only supplementary notes.
\usepackage[T1]{fontenc}
\usepackage{amsmath,amssymb,amsthm}
\usepackage{hyperref}
\usepackage{cleveref}
\DeclareMathOperator{\Map}{Map}
\DeclareMathOperator{\Fun}{Fun}
\newcommand{\iso}{\xrightarrow{\sim}}
\newtheorem{lemma}{Supplementary Lemma}
\newtheorem{corollary}[lemma]{Supplementary Corollary}
\crefname{lemma}{Supplementary Lemma}{Supplementary Lemmas}
\Crefname{lemma}{Supplementary Lemma}{Supplementary Lemmas}
\crefname{corollary}{Supplementary Corollary}{Supplementary Corollaries}
\Crefname{corollary}{Supplementary Corollary}{Supplementary Corollaries}

\begin{document}
\section*{Inverting the walking morphism}

The \href{geometric-realizations.html}{geometric realization of $[1]$}
is contractible. The checked proof identifies the functors out of $[1]$
that invert every arrow by evaluating their inversion data at the identity
functor $[1]\to[1]$. This is an alternative to the manuscript's
argument enumerating the three arrows of $[1]$.

Use the Chapter 1 structures, walking-morphism and endpoint data,
the Segal and commutative-square axioms, and Rezk.
For a category $D$, write
\[
 V_D=\Map^{\Map([1],[1])}([1],D),\qquad
 u\colon V_D\hookrightarrow\Map([1],D).
\]
Thus $V_D$ consists of functors inverting all arrows of $[1]$.
Let $c\colon D^\simeq\to\Map([1],D)$ be the constant-arrow functor.

\begin{lemma}
\label[lemma]{lem:Supplementary_Interval_Constants}
There is a functor $k\colon D^\simeq\to V_D$ with $uk=c$.
\end{lemma}
\begin{proof}[Proof idea]
Constant arrows factor through the terminal category. That category is a
groupoid, and a functor out of a groupoid inverts every arrow. Apply this
to the universal family of objects of $D$, then restrict the inverting
family along $[1]\to\star$. The restriction computation identifies its
underlying family with $c$.
\end{proof}

\begin{lemma}
\label[lemma]{lem:Supplementary_Inverted_Identity}
There is a functor $o\colon V_D\to D^\simeq$ with $co=u$.
The functor $k$ of \Cref{lem:Supplementary_Interval_Constants} is an
equivalence with inverse $o$.
\end{lemma}
\begin{proof}[Proof idea]
The defining pullback of $V_D$ supplies a family of isomorphisms indexed
by the arrows of $[1]$. Evaluate this family at the point of
$\Map([1],[1])$ naming the identity functor. Its underlying arrow is the
original functor $[1]\to D$: this is the identity computation for the
mapping action, with the naming comparisons included.

Rezk identifies the anima of isomorphisms in $D$ with $D^\simeq$, by
sending an object to its identity. Lifting the evaluated family through
this equivalence gives $o$ and the comparison $co=u$.
Both $u$ and $c$ are embeddings; for $c$ this follows from Rezk and the
embedding of isomorphisms among all arrows. The equations $uk=c$ and
$co=u$ therefore imply $ok=\operatorname{id}$ and
$ko=\operatorname{id}$ by reflection. These are identifications, not
strict equalities of the chosen functors.
\end{proof}

\begin{corollary}
\label[corollary]{cor:Supplementary_Interval_Localization}
In the setting of the localization definition, let $l\colon[1]\to T$
be a localization at all arrows of $[1]$. Then $T$ is contractible.
\end{corollary}
\begin{proof}[Proof idea]
For every $D$, the localization property gives an equivalence
$\rho_D\colon\Map(T,D)\iso V_D$.
The composite of restriction along $T\to\star$ with $\rho_D$ is
identified with $k$. To check this comparison, compose with the embedding
$u$: both routes become the constant-arrow family, using the unit and
composition comparisons for restriction and the terminal comparison
$(T\to\star)l=([1]\to\star)$. Reflect through $u$.

Since $k$ and $\rho_D$ are equivalences, restriction along $T\to\star$
is an equivalence for every $D$. Detection by mapping animae shows that
$T\to\star$ is an equivalence. No assumption that $T$ is already a
groupoid is used.
\end{proof}
\end{document}
